% \documentclass[review=false, screen=true]{acmart}
\documentclass[format=acmsmall, review=false, screen=true]{acmart}

% Package to generate and customize Algorithm as per ACM style
\usepackage[linesnumbered,ruled,vlined]{algorithm2e}
\renewcommand{\algorithmcfname}{ALGORITHM}

\usepackage{booktabs}
\usepackage{tabularx}

\usepackage{amsmath,amssymb,amsfonts}

\usepackage{enumitem}

\SetKwRepeat{Do}{do}{while}

\usepackage{lmodern}
\usepackage{courier}

\usepackage{listings}
\lstset{
  language=Java,
  backgroundcolor=\color{white},
  basicstyle=\ttfamily\small,
  breakatwhitespace=false,
  breaklines=true,
  captionpos=b,
  commentstyle=\color{red},
  frame=single,
  keepspaces=true,
  keywordstyle=\color{blue},
  showspaces=false,
  showstringspaces=false,
  showtabs=false,
  stringstyle=\color{brown},
  tabsize=2,
}

\usepackage{setspace}



% Document starts
\begin{document}
% Title portion. Note the short title for running heads
\title[COSC 3P91 -- Assignment 2]{COSC 3P91 -- Assignment 2 -- Village War Strategy Game}

\author{Group Member 1}
\affiliation{%
  \institution{Brock University}
  \streetaddress{1812 Sir Isaac Brock Way}
  \city{St. Catharines}
  \state{ON}
  \postcode{L2S 3A1}
  \country{Canada}
}

\begin{abstract}
This document describes the design and implementation of a console-based village
war strategy game developed for COSC 3P91 Assignment 2.  The game allows a human
player to build and upgrade a village, train an army, explore and attack
randomly-generated NPC villages, and accumulate loot to rank up on a
leaderboard.  The implementation demonstrates several advanced Java
object-oriented programming concepts mandated by the assignment specification,
including generics, local and anonymous classes, lambda expressions and method
references, custom checked exceptions, Java standard utility collections, and
the Stream API for I/O-oriented processing.
\end{abstract}

\maketitle

% ------------------------------------------------------------------
\section{Java Application Design and Implementation}
% ------------------------------------------------------------------

The game is implemented as a multi-package Java application that follows the
UML class diagram supplied with the assignment.  Five packages organise the
source tree:

\begin{itemize}[noitemsep]
  \item \texttt{game} -- core game logic: \texttt{GameEngine}, \texttt{Village},
        \texttt{Army}, \texttt{Player}, \texttt{Defence}, \texttt{CollectResources},
        \texttt{WallClock}, and the generic \texttt{Repository<T>}.
  \item \texttt{gameelements} -- buildings, habitants, and resources:
        abstract hierarchies \texttt{Building}, \texttt{Habitant}/\texttt{Fighter}/\texttt{Peasant},
        and \texttt{Resource}, together with all concrete sub-classes.
  \item \texttt{utility} -- combat support: \texttt{Arbitrer}, \texttt{AttackOutcome},
        \texttt{GameMap}, \texttt{Position}, and \texttt{Region}.
  \item \texttt{exceptions} -- four custom checked exceptions.
  \item \texttt{gui} -- \texttt{GraphicalInterface}, the console-based UI renderer.
\end{itemize}

The application entry point is \texttt{Main.java}, which hosts the game loop
and dispatches all player menu choices.

% ------------------------------------------------------------------
\subsection{Building}
% ------------------------------------------------------------------

A player selects a building type from a menu in \texttt{Main.buildBuilding()}.
The cost constants are stored as \texttt{public static final double} fields on
each concrete building class (e.g.\ \texttt{Farm.COST\_GOLD}).  The actual
construction is performed by \texttt{Village.build(Building)}, which:
\begin{enumerate}[noitemsep]
  \item Throws \texttt{BuildingLimitExceededException} if the village already
        holds the maximum of 20 buildings.
  \item Calls \texttt{Village.checkResources()} and throws
        \texttt{InsufficientResourcesException} if any resource is short.
  \item Deducts the resources and appends the building to the internal
        \texttt{List<Building>}.
\end{enumerate}

% ------------------------------------------------------------------
\subsection{Training}
% ------------------------------------------------------------------

Military and worker units are trained through \texttt{Main.trainUnit()}.  The
player selects a unit type (Soldier, Archer, Knight, Catapult, Gold Miner,
Iron Miner, or Lumberman).  A new instance of the chosen concrete
\texttt{Habitant} sub-class is created, resources are deducted by
\texttt{Village.train()}, and the unit is appended to the village's
\texttt{List<Habitant>}.  Fighters are simultaneously added to
\texttt{Player.getArmy()} via \texttt{Army.addFighter()}, which throws
\texttt{InvalidOperationException} when the 10-unit army cap is reached.

% ------------------------------------------------------------------
\subsection{Upgrading}
% ------------------------------------------------------------------

\texttt{Village.upgrade(Building)} upgrades any building already present in the
village.  It enforces two constraints:
\begin{itemize}[noitemsep]
  \item The building's level must be below the \texttt{VillageHall} level (the
        hall therefore acts as a global upgrade gate).
  \item Resources must cover the upgrade cost returned by
        \texttt{getUpgradeCostGold/Iron/Lumber()}.
\end{itemize}
After passing all checks, \texttt{Building.upgrade()} increments the level,
adds 50 hit-points, and calls the abstract hook \texttt{applyUpgradeBonus()},
which each subclass overrides to improve its own statistics (e.g.\
\texttt{ArcherTower} increases its damage; \texttt{Farm} increases its food
capacity).

% ------------------------------------------------------------------
\subsection{Generating Villages}
% ------------------------------------------------------------------

\texttt{GameEngine.generateNpcVillages(int count)} creates \textit{count} NPC
villages and stores them in \texttt{Repository<Village> npcVillages}.  The
method uses a \textbf{local class} \texttt{NpcBuilder} to encapsulate village
construction logic, and scales the NPC's buildings and resources based on a
randomly selected difficulty value (1--3).  Village names are drawn from a
pre-defined array of fantasy names.

% ------------------------------------------------------------------
\subsection{Attack Exploring}
% ------------------------------------------------------------------

\texttt{GameEngine.getAvailableTargets()} returns all NPC villages sorted by
defence score in ascending order, using an \textbf{anonymous} \texttt{Comparator}
class.  The player sees the sorted list in \texttt{Main.exploreTargets()} and
can pick a target by index.  The selected \texttt{Village} is stored for the
subsequent attack action.

% ------------------------------------------------------------------
\subsection{Attacking}
% ------------------------------------------------------------------

\texttt{Main.attackVillage()} calls \texttt{Army.attack(Village defender)},
which delegates to \texttt{Arbitrer.judgeAttack()}.  The arbitration formula
is:
\[
  \text{effectiveAttack} = \text{attackScore} + \mathcal{U}(0, 50)
\]
where \(\text{attackScore} = \text{totalDamage} \times \sqrt{|\text{fighters}|}\)
to reward larger squads.  If \texttt{effectiveAttack} exceeds the
defender's defence score the attack succeeds; loot is proportional to the
victory margin and capped at 30\,\% of the defender's stores.  On failure,
roughly 30\,\% of the attacking army is removed at random (casualties).  The
outcome is encapsulated in \texttt{AttackOutcome} and the loot is credited to
the player's village via \texttt{Village.addResources()}.

% ------------------------------------------------------------------
\subsection{Game Engine}
% ------------------------------------------------------------------

\texttt{GameEngine} is the central controller.  It owns the \texttt{WallClock}
(a tick-based time simulation), the player and NPC village repositories, and
the \texttt{GameMap} (a 1000$\times$1000 grid of named \texttt{Position}
entries).  Each iteration of the main loop calls \texttt{engine.tick()} to
advance game time.  The engine also maintains the leaderboard: players are
ranked by score using \texttt{Comparator.comparingInt(Player::getScore).reversed()},
a method reference passed to the Stream \texttt{sorted()} terminal operation.

\begin{table}[h]
  \centering
  \caption{Building types, starting hit-points, and upgrade costs.}
  \begin{tabular}{lrrrr}
    \toprule
    \textbf{Building} & \textbf{HP} & \textbf{Gold} & \textbf{Iron} & \textbf{Lumber} \\
    \midrule
    Village Hall   & 500 & 100 & 50 & 50 \\
    Farm           & 200 & 25  & 15 & 10 \\
    Gold Mine      & 300 & 50  & 25 & 25 \\
    Iron Mine      & 300 & 40  & 30 & 20 \\
    Lumber Mill    & 300 & 15  & 5  & 30 \\
    Archer Tower   & 350 & 30  & 20 & 15 \\
    Cannon         & 400 & 50  & 40 & 20 \\
    \bottomrule
  \end{tabular}
\end{table}

% ------------------------------------------------------------------
\section{Object-Oriented Principles and Java Concepts}
% ------------------------------------------------------------------

% ------------------------------------------------------------------
\subsection{Generics}
% ------------------------------------------------------------------

The class \texttt{Repository<T>} (Listing~\ref{lst:repository}) is the
canonical demonstration of generics.  It is instantiated twice in
\texttt{GameEngine}: once as \texttt{Repository<Player>} and once as
\texttt{Repository<Village>}, giving strongly-typed access to each store
without raw-type casts.  The backing data structure is a \texttt{TreeMap<String, T>},
preserving lexicographic key order.

The method \texttt{Army.addAll(List<? extends Fighter>)} demonstrates an
upper-bounded wildcard: any list whose element type is \texttt{Fighter} or a
subtype is accepted, enabling reuse with heterogeneous fighter lists without
sacrificing type safety.

\begin{spacing}{0.8}
\begin{lstlisting}[language=Java, caption={Generic Repository (excerpt).}, label=lst:repository]
public class Repository<T> {
    private final TreeMap<String, T> store = new TreeMap<>();

    public void put(String key, T entity) { store.put(key, entity); }

    public List<T> findWhere(Predicate<T> predicate) {
        return store.values().stream()
                .filter(predicate)
                .collect(Collectors.toList());
    }
    public void forEach(Consumer<T> action) {
        store.values().forEach(action);
    }
}
\end{lstlisting}
\end{spacing}

% ------------------------------------------------------------------
\subsection{Local and Anonymous Classes}
% ------------------------------------------------------------------

\textbf{Local class.}~\texttt{GameEngine.generateNpcVillages()} declares a local
class \texttt{NpcBuilder} whose sole \texttt{build()} method constructs and
populates an NPC village.  Scoping the class locally keeps the builder logic
cohesive and prevents it from polluting the package namespace.

A second local class \texttt{BuildingMerger} appears inside
\texttt{Defence.getAllDefenceBuildings()}: it merges the archer-tower and
cannon lists into a single \texttt{List<Building>} using an upper-bounded
wildcard parameter.

\textbf{Anonymous class.}~\texttt{GameEngine.getAvailableTargets()} passes an
anonymous \texttt{Comparator<Village>} to \texttt{List.sort()} for ordering
targets by ascending defence score.  Similarly,
\texttt{Village.getBuildingsSortedByLevel()} uses an anonymous
\texttt{Comparator<Building>} to sort by descending level.

\begin{spacing}{0.8}
\begin{lstlisting}[language=Java, caption={Local class NpcBuilder and anonymous Comparator.}, label=lst:localanon]
// Local class inside generateNpcVillages()
class NpcBuilder {
    Village build(String name, int difficulty) {
        Village v = new Village(name, 200 + difficulty * 100, ...);
        if (difficulty >= 1) v.getBuildingsMutable().add(new ArcherTower());
        ...
        return v;
    }
}

// Anonymous Comparator inside getAvailableTargets()
targets.sort(new Comparator<Village>() {
    @Override
    public int compare(Village a, Village b) {
        return Double.compare(a.getDefenceScore(), b.getDefenceScore());
    }
});
\end{lstlisting}
\end{spacing}

% ------------------------------------------------------------------
\subsection{Lambda Expressions and Method References}
% ------------------------------------------------------------------

Lambdas and method references appear throughout the codebase:

\begin{itemize}[noitemsep]
  \item \texttt{Village.getDefenceScore()} -- \texttt{mapToDouble(b -> b.getLevel() * b.getHitPoints()).sum()}.
  \item \texttt{Village.getAttackScore()} -- stream pipeline with a filter lambda and a \texttt{mapToDouble} lambda.
  \item \texttt{Village.getStatusSummary()} -- \texttt{buildings.forEach(b -> sb.append(...))} and \texttt{habitants.forEach(h -> ...)}.
  \item \texttt{Army.recalculateDamage()} -- \texttt{mapToDouble(Fighter::damage).sum()}, using a method reference.
  \item \texttt{Army.getSummary()} -- \texttt{Collectors.groupingBy(f -> f.getClass().getSimpleName(), Collectors.counting())} followed by \texttt{forEach((type, count) -> ...)}.
  \item \texttt{GameEngine.getLeaderboard()} -- \texttt{sorted(Comparator.comparingInt(Player::getScore).reversed())}, a method reference composed with a reversal combinator.
  \item \texttt{CollectResources.collect()} -- \texttt{map(Peasant::work)} and subsequent \texttt{filter}/\texttt{mapToDouble} pipelines for each building type.
  \item \texttt{Defence.syncWithVillage()} -- \texttt{forEach(archerTowers::add)} and \texttt{forEach(cannons::add)}, instance method references.
\end{itemize}

\begin{spacing}{0.8}
\begin{lstlisting}[language=Java, caption={Lambda and method reference examples.}, label=lst:lambda]
// Village.getDefenceScore()
return buildings.stream()
        .mapToDouble(b -> b.getLevel() * b.getHitPoints())
        .sum();

// Army.recalculateDamage() -- method reference
damage = fighters.stream()
        .mapToDouble(Fighter::damage)
        .sum();

// GameEngine.getLeaderboard() -- method reference + combinator
return playerRepository.getAll().stream()
        .sorted(Comparator.comparingInt(Player::getScore).reversed())
        .collect(Collectors.toList());
\end{lstlisting}
\end{spacing}

% ------------------------------------------------------------------
\subsection{Custom Exceptions}
% ------------------------------------------------------------------

Four checked exceptions in the \texttt{exceptions} package model domain-specific
error conditions:

\begin{itemize}[noitemsep]
  \item \texttt{InsufficientResourcesException} -- thrown by
        \texttt{Village.checkResources()} when gold, iron, or lumber is short.
        Stores the resource type, required amount, and available amount.
  \item \texttt{BuildingLimitExceededException} -- thrown by \texttt{Village.build()}
        when the 20-building cap is reached.  Stores current count and cap.
  \item \texttt{MaxLevelReachedException} -- thrown by \texttt{Building.upgrade()}
        when the level-5 cap is reached, or by \texttt{Village.upgrade()} when
        the VillageHall gate is hit.
  \item \texttt{InvalidOperationException} -- a general-purpose exception for
        illegal state transitions (e.g.\ attacking with an empty army,
        removing a non-existent building).
\end{itemize}

Each exception class provides two constructors: a single-string constructor for
programmatic use and a structured constructor that generates a descriptive message
from the relevant parameters.  The \texttt{Main} game loop catches all four
exceptions and displays a user-friendly error line prefixed with \texttt{[!]}.

\begin{spacing}{0.8}
\begin{lstlisting}[language=Java, caption={InsufficientResourcesException.}, label=lst:exception]
public class InsufficientResourcesException extends Exception {
    private final String resourceType;
    private final double required;
    private final double available;

    public InsufficientResourcesException(
            String resourceType, double required, double available) {
        super(String.format(
            "Insufficient %s: required %.1f but only %.1f available.",
            resourceType, required, available));
        ...
    }
}
\end{lstlisting}
\end{spacing}

% ------------------------------------------------------------------
\subsection{Java Utility Classes (Collections)}
% ------------------------------------------------------------------

The following standard collection types are used:

\begin{itemize}[noitemsep]
  \item \texttt{ArrayList<Building>}, \texttt{ArrayList<Habitant>},
        \texttt{ArrayList<Fighter>} -- the primary mutable collections inside
        \texttt{Village} and \texttt{Army}.
  \item \texttt{TreeMap<String, T>} -- the backing store of \texttt{Repository<T>},
        chosen for its guaranteed sorted key order and $O(\log n)$ access.
  \item \texttt{Collections.unmodifiableList()} -- used by \texttt{Village.getBuildings()}
        and \texttt{Village.getHabitants()} to expose read-only views of internal
        state, protecting encapsulation.
  \item \texttt{Map<String, Long>} produced by
        \texttt{Collectors.groupingBy(..., Collectors.counting())} inside
        \texttt{Army.getSummary()} to count fighters by type.
\end{itemize}

% ------------------------------------------------------------------
\subsection{I/O and Streams}
% ------------------------------------------------------------------

The Stream API drives much of the game's data processing.  Key uses:

\begin{itemize}[noitemsep]
  \item \textbf{Filtering:}  \texttt{Village.getFighters()} and
        \texttt{Village.getPeasants()} filter the \texttt{habitants} list by
        \texttt{instanceof} predicates.
  \item \textbf{Mapping:}  \texttt{CollectResources.collect()} calls
        \texttt{map(Peasant::work)} to transform each peasant into the
        \texttt{Resource} it produces.
  \item \textbf{Aggregation (reduce / sum):}
        \texttt{Defence.getTotalDamage()} accumulates archer-tower and cannon
        damage via \texttt{mapToDouble(...).reduce(0, Double::sum)}.
  \item \textbf{Grouping:}  \texttt{Army.getSummary()} groups fighters by
        class name and counts them with \texttt{Collectors.groupingBy}.
  \item \textbf{Sorting:}  \texttt{Village.getFightersSortedByDamage()} returns
        fighters ordered by descending damage using a composed comparator:
        \texttt{Comparator.comparingDouble(Fighter::damage).reversed()}.
\end{itemize}

Console I/O uses a \texttt{Scanner} wrapping a \texttt{BufferedReader} around
\texttt{System.in} in \texttt{Main.main()}, giving buffered line-oriented input
with minimal latency.

\begin{spacing}{0.8}
\begin{lstlisting}[language=Java, caption={Stream-based resource collection in CollectResources.}, label=lst:streams]
List<Resource> produced = peasants.stream()
        .map(Peasant::work)
        .collect(Collectors.toList());

totalGold += buildings.stream()
        .filter(b -> b instanceof GoldMine)
        .mapToDouble(b -> ((GoldMine) b).getGoldProd())
        .sum();
\end{lstlisting}
\end{spacing}

% ------------------------------------------------------------------
\subsection{Performance Analysis}
% ------------------------------------------------------------------

The game runs in a tight console loop.  Each iteration is $O(B + H)$ where $B$
is the number of buildings and $H$ the number of habitants, both bounded by
constants ($B \le 20$, $H \le 10$ fighters plus unlimited peasants in
practice).  The \texttt{Repository<T>} uses a \texttt{TreeMap} giving
$O(\log n)$ key-based lookup; for the small NPC village set ($n = 5$ by default)
this is negligible.  Combat resolution in \texttt{Arbitrer} is $O(1)$.

% ------------------------------------------------------------------
\subsection{Conclusions}
% ------------------------------------------------------------------

The implementation covers all six required behaviours (Building, Training,
Upgrading, Generating Village, Attack Exploring, and Attacking) and
demonstrates each of the six mandated Java concepts.  The class hierarchy is
cleanly separated along lines of responsibility: \texttt{game} classes
orchestrate high-level behaviour, \texttt{gameelements} classes model domain
entities, \texttt{utility} classes handle combat mathematics, and
\texttt{exceptions} classes communicate error conditions in a type-safe way.
The design can be extended -- for example by adding new unit types, a
network multiplayer layer, or a graphical front-end -- without modifying
existing classes.


\bibliographystyle{ACM-Reference-Format-Journals}
\bibliography{acmsmall-sample-bibfile}


\appendix

% ------------------------------------------------------------------
\section{Main.java -- Game Loop (excerpt)}
% ------------------------------------------------------------------
\begin{spacing}{0.8}
\begin{lstlisting}[language=Java]
public static void main(String[] args) {
    scanner   = new Scanner(new BufferedReader(
                    new InputStreamReader(System.in)));
    engine    = new GameEngine();
    player    = new Player(1, villageName);
    engine.registerPlayer(player);
    engine.generateNpcVillages(5);
    engine.setRunning(true);

    while (engine.isRunning()) {
        engine.tick();
        ui.renderMenu();
        String input = scanner.nextLine().trim();
        try {
            handleMenuChoice(input);
        } catch (InsufficientResourcesException e) {
            System.out.println("  [!] " + e.getMessage());
        } catch (MaxLevelReachedException e) {
            System.out.println("  [!] " + e.getMessage());
        } catch (BuildingLimitExceededException e) {
            System.out.println("  [!] " + e.getMessage());
        } catch (InvalidOperationException e) {
            System.out.println("  [!] " + e.getMessage());
        }
    }
}
\end{lstlisting}
\end{spacing}

\clearpage

% ------------------------------------------------------------------
\section{Village.java -- Score Calculations (excerpt)}
% ------------------------------------------------------------------
\begin{spacing}{0.8}
\begin{lstlisting}[language=Java]
// Defence score: sum of level * hitPoints over all buildings
public double getDefenceScore() {
    return buildings.stream()
            .mapToDouble(b -> b.getLevel() * b.getHitPoints())
            .sum();
}

// Attack score: sum of damage * level for each Fighter habitant
public double getAttackScore() {
    return habitants.stream()
            .filter(h -> h instanceof Fighter)
            .map(h -> (Fighter) h)
            .mapToDouble(f -> f.damage() * f.getLevel())
            .sum();
}

// Buildings sorted by level (anonymous Comparator)
public List<Building> getBuildingsSortedByLevel() {
    List<Building> sorted = new ArrayList<>(buildings);
    Collections.sort(sorted, new Comparator<Building>() {
        @Override
        public int compare(Building a, Building b) {
            return Integer.compare(b.getLevel(), a.getLevel());
        }
    });
    return sorted;
}
\end{lstlisting}
\end{spacing}

\clearpage

% ------------------------------------------------------------------
\section{Army.java -- Attack and Wildcard (excerpt)}
% ------------------------------------------------------------------
\begin{spacing}{0.8}
\begin{lstlisting}[language=Java]
// Attack score with squad-size synergy
public double getAttackScore() {
    if (fighters.isEmpty()) return 0;
    return damage * Math.sqrt(fighters.size());
}

// Bounded wildcard: accepts any List<? extends Fighter>
public void addAll(List<? extends Fighter> newFighters)
        throws InvalidOperationException {
    for (Fighter f : newFighters) addFighter(f);
}

// Method reference: Fighter::damage
private void recalculateDamage() {
    damage = fighters.stream()
            .mapToDouble(Fighter::damage)
            .sum();
}
\end{lstlisting}
\end{spacing}

\end{document}
